% Vietnamese translation for OI Public Library
% Source-PDF: IOI中国国家候选队论文/2001/张一飞.pdf
\documentclass[11pt]{article}
\usepackage{oipl-vn}

\title{Thuật toán độ chính xác cao để tính \(N!\)}
\author{Zhang Yifei}
\date{}

\begin{document}
\maketitle

\origtitle{High-precision algorithms for computing \(N!\)}
\sourcepath{IOI China candidate team papers / 2001 / Zhang Yifei.pdf}

\tableofcontents

\section{Dẫn nhập}

Các thuật toán trong bài này chủ yếu nhắm tới ngôn ngữ Pascal. Bài viết đặt
ra ba câu hỏi: ta hiểu gì về số học độ chính xác cao, ta đã từng dùng những
cấu trúc dữ liệu nào cho số nguyên lớn, và ta đã suy nghĩ kỹ về cách tối ưu
thuật toán hay chưa. Với bài toán tính \(N!\), ta sẽ thấy một số cách cải
thiện có thể làm chương trình nhanh hơn nhiều lần.

\section{Về độ chính xác cao}

\subsection{Các kiểu số nguyên chuẩn trong Pascal}

Pascal cung cấp một số kiểu số nguyên chuẩn với miền giá trị cố định:
\[
\begin{array}{lll}
\text{Kiểu dữ liệu} & \text{Miền giá trị} \\
\hline
\texttt{Shortint} & -128 & \text{đến }127 \\
\texttt{Byte}     & 0    & \text{đến }255 \\
\texttt{Integer}  & -32768 & \text{đến }32767 \\
\texttt{Word}     & 0 & \text{đến }65535 \\
\texttt{Longint}  & -2147483648 & \text{đến }2147483647 \\
\texttt{Comp}     & -9.2\cdot 10^{18} & \text{đến }9.2\cdot 10^{18}.
\end{array}
\]
Tuy \texttt{Comp} thuộc nhóm kiểu thực, trên thực tế nó biểu diễn một số
nguyên 64 bit.

\subsection{Tư tưởng cơ bản của thuật toán độ chính xác cao}

Các kiểu số nguyên chuẩn trong Pascal nhiều nhất chỉ xử lý trực tiếp được
những số nguyên nằm trong khoảng xấp xỉ từ \(-2^{63}\) đến \(2^{63}\). Khi
cần tính với số nguyên lớn hơn, ta phải dùng thuật toán độ chính xác cao.

Tư tưởng cơ bản là chia số nguyên lớn, vốn không thể xử lý trực tiếp, thành
nhiều đoạn số nguyên nhỏ có thể xử lý trực tiếp. Khi đó phép toán trên số lớn
được chuyển thành các phép toán trên những đoạn nhỏ ấy.

\section{Lựa chọn cấu trúc dữ liệu}

\subsection{Giữ càng nhiều chữ số trong mỗi đoạn càng tốt}

Xét ví dụ tính tổng của hai số có 15 chữ số. Có ba cách chia số:
\begin{itemize}
  \item Cách thứ nhất: chia thành 15 đoạn nhỏ, mỗi đoạn có 1 chữ số. Ta cần
  15 phép cộng một chữ số.
  \item Cách thứ hai: chia thành 5 đoạn nhỏ, mỗi đoạn có 3 chữ số. Ta cần 5
  phép cộng ba chữ số.
  \item Cách thứ ba: kiểu \texttt{Comp} có thể xử lý trực tiếp số nguyên 15
  chữ số, nên chỉ cần 1 phép cộng.
\end{itemize}

Khi dùng \texttt{Integer}, cộng hai số một chữ số và cộng hai số ba chữ số có
tốc độ như nhau, vì vậy cách thứ hai hiệu quả hơn cách thứ nhất. Dù thao tác
trên \texttt{Comp} chậm hơn thao tác trên \texttt{Integer}, số lần cộng giảm
rất mạnh; thực nghiệm cho thấy cách thứ ba vẫn nhanh hơn cách thứ hai.

\subsection{Dùng kiểu \texttt{Comp}}

Trong tính toán độ chính xác cao, mỗi đoạn nhỏ có thể được biểu diễn bằng
kiểu \texttt{Comp}. Kiểu này có khoảng 19 đến 20 chữ số có nghĩa. Tùy phép
toán, ta chọn số chữ số lưu trong mỗi đoạn như sau:
\begin{itemize}
  \item Khi cộng hai số độ chính xác cao, mỗi đoạn có thể giữ 17 chữ số.
  \item Khi nhân một số độ chính xác cao với một số nguyên không quá \(m\)
  chữ số, mỗi đoạn có thể giữ \(18-m\) chữ số.
  \item Khi nhân hai số độ chính xác cao, mỗi đoạn có thể giữ 9 chữ số.
\end{itemize}

Nếu mỗi đoạn nhỏ giữ \(k\) chữ số thập phân, ta có thể xem đoạn đó là một chữ
số trong hệ cơ số \(10^k\). Trong bài này, ta gọi cách biểu diễn ấy là
\term{số cơ số cao}, và gọi một chữ số cơ số cao là một \term{số đơn chính
xác}.

\subsection{Biểu diễn nhị phân}

Dùng biểu diễn nhị phân có thể cải thiện hiệu quả thời gian và không gian
trong quá trình tính toán. Tuy nhiên, đề bài thường yêu cầu in kết quả ở hệ
thập phân, nên cần thêm một bước đổi cơ số rất tốn thời gian. Vì thế phương
pháp này thường không được dùng trong thi đấu, và bài viết không đi sâu vào
cách làm đó.

\section{Tối ưu hóa thuật toán}

Các phần sau phân tích độ phức tạp của phép nhân độ chính xác cao, phép nhân
liên tiếp, kỹ thuật đặt bộ đệm, cách tính giai thừa bằng phân tích thừa số
nguyên tố, phương pháp lũy thừa nhị phân, và bước điều chỉnh sau khi phân
tích thừa số nguyên tố.

\subsection{Độ phức tạp của phép nhân độ chính xác cao}

Khi nhân một số có \(n\) chữ số cơ số cao với một số có \(m\) chữ số cơ số
cao, thuật toán nhân thông thường cần \(n m\) phép nhân giữa các chữ số cơ số
cao. Kết quả có thể có \(n+m-1\) hoặc \(n+m\) chữ số cơ số cao.

\subsection{Độ phức tạp của phép nhân liên tiếp}

Xét ví dụ tính
\[
  5\cdot 6\cdot 7\cdot 8.
\]
Nếu nhân theo thứ tự từ trái sang phải, ta có
\[
\begin{array}{rcll}
5\cdot 6 &=& 30,   & 1\cdot 1=1 \text{ phép nhân},\\
30\cdot 7 &=& 210, & 2\cdot 1=2 \text{ phép nhân},\\
210\cdot 8 &=& 1680, & 3\cdot 1=3 \text{ phép nhân}.
\end{array}
\]
Tổng cộng có 6 phép nhân. Nếu nhóm khác đi,
\[
\begin{array}{rcll}
5\cdot 6 &=& 30, & 1\cdot 1=1 \text{ phép nhân},\\
7\cdot 8 &=& 56, & 1\cdot 1=1 \text{ phép nhân},\\
30\cdot 56 &=& 1680, & 2\cdot 2=4 \text{ phép nhân},
\end{array}
\]
tổng số phép nhân vẫn là 6.

Đặc điểm của ví dụ này là khi một số \(n\) chữ số nhân với một số \(m\) chữ
số, kết quả có đúng \(n+m\) chữ số. Nếu giả thiết ấy luôn đúng, thì với \(n\)
số đơn chính xác, bất kể chọn thứ tự nhân như thế nào, số phép nhân luôn là
\[
  \frac{n(n-1)}{2}.
\]

Chứng minh bằng quy nạp như sau. Gọi \(F(n)\) là số phép nhân cần dùng. Ta có
\(F(1)=0\). Giả sử với mọi \(k<n\), \(F(k)=k(k-1)/2\). Ở lần nhân cuối, ta
nhân một số \(k\) chữ số với một số \(n-k\) chữ số, nên
\[
\begin{aligned}
F(n)
&=F(k)+F(n-k)+k(n-k)\\
&=\frac{k(k-1)}{2}
  +\frac{(n-k)(n-k-1)}{2}
  +k(n-k)\\
&=\frac{n(n-1)}{2}.
\end{aligned}
\]
Giá trị này không phụ thuộc vào cách chọn \(k\).

\subsection{Đặt bộ đệm}

Trường hợp kết quả của phép nhân có thể chỉ có \(n+m-1\) chữ số cho thấy thứ
tự nhân bắt đầu có ảnh hưởng. Ví dụ tính
\[
  9\cdot 8\cdot 3\cdot 2.
\]
Nếu nhân theo thứ tự từ trái sang phải:
\[
\begin{array}{rcll}
9\cdot 8 &=& 72, & 1\cdot 1=1 \text{ phép nhân},\\
72\cdot 3 &=& 216, & 2\cdot 1=2 \text{ phép nhân},\\
216\cdot 2 &=& 432, & 3\cdot 1=3 \text{ phép nhân}.
\end{array}
\]
Tổng cộng có 6 phép nhân. Nếu trước hết tính \(9\cdot 8\) và \(3\cdot 2\),
rồi mới nhân hai kết quả:
\[
\begin{array}{rcll}
9\cdot 8 &=& 72, & 1\cdot 1=1 \text{ phép nhân},\\
3\cdot 2 &=& 6, & 1\cdot 1=1 \text{ phép nhân},\\
72\cdot 6 &=& 432, & 2\cdot 1=2 \text{ phép nhân},
\end{array}
\]
tổng cộng chỉ còn 4 phép nhân.

Ta xét tích của \(k+t\) số đơn chính xác
\[
  a_1a_2\cdots a_k a_{k+1}\cdots a_{k+t}.
\]
Giả sử \(a_1a_2\cdots a_k\) đã được tính và kết quả là một số có \(m\) chữ số
cơ số cao; đồng thời \(a_{k+1}\cdots a_{k+t}\) vẫn là một số đơn chính xác.
Nếu tiếp tục nhân tuần tự, ta cần \(t\) phép nhân kiểu ``số \(m\) chữ số nhân
với số 1 chữ số'', tức khoảng \(mt\) phép nhân giữa chữ số cơ số cao. Nếu
trước hết tính \(a_{k+1}\cdots a_{k+t}\), rồi nhân một lần với kết quả trước,
ta chỉ cần khoảng \(m+t\) phép nhân.

Trong cách cài đặt có bộ đệm, ta gom nhiều số đơn chính xác lại cho đến khi
tích của chúng sắp vượt quá khả năng chứa của kiểu chuẩn, rồi mới nhân tích
đó vào số độ chính xác cao. Khi tính tích của \(m\) số đơn chính xác và kết
quả cuối cùng có \(n\) chữ số cơ số cao, số phép nhân xấp xỉ
\[
  \frac{n(n-1)}{2},
\]
gần như không phụ thuộc trực tiếp vào \(m\). Có thể hình dung ta chia \(m\)
thừa số thành \(n\) nhóm, mỗi nhóm có tích vẫn là một số đơn chính xác; phần
tốn thời gian chính là nhân \(n\) số thu được. Lúc đó mô hình
``\(n\) chữ số nhân \(m\) chữ số cho \(n+m\) chữ số'' gần đúng, nên công thức
trên áp dụng được.

\subsection{Kích thước bộ đệm}

Giả sử kiểu dữ liệu chuẩn được chọn có thể xử lý trực tiếp tối đa \(t\) chữ
số thập phân. Đặt kích thước bộ đệm là \(k\) chữ số thập phân; khi đó mỗi
đoạn nhỏ của số độ chính xác cao giữ \(t-k\) chữ số thập phân. Nếu kết quả
cuối cùng có \(n\) chữ số thập phân, thì số phép nhân xấp xỉ
\[
  \frac{k}{n-k}\sum_{i=1}^{n/k} i
  =\frac{(n+k)n}{2k(t-k)},
\]
trong đó \(k\) nhỏ hơn \(n\) rất nhiều.

Để số phép nhân nhỏ nhất, ta cần làm cho \(k(t-k)\) lớn nhất. Giá trị tối ưu
đạt được khi \(k=t/2\). Vì vậy, hiệu quả cao nhất xuất hiện khi kích thước bộ
đệm bằng kích thước mỗi đoạn nhỏ. Trong phép nhân liên tiếp thông thường, có
thể dùng \texttt{Comp} làm kiểu cơ sở, mỗi đoạn nhỏ giữ 9 chữ số thập phân,
và bộ đệm cũng giữ 9 chữ số thập phân.

\section{Tính giai thừa bằng phân tích thừa số nguyên tố}

Ví dụ:
\[
  10! = 2^8\cdot 3^4\cdot 5^2\cdot 7.
\]
Độ phức tạp của việc phân tích \(n!\) thành thừa số nguyên tố nhỏ hơn rất
nhiều so với \(n\log n\), nên có thể xem là không đáng kể so với phần nhân độ
chính xác cao.

So với thuật toán thông thường, sau khi phân tích thừa số nguyên tố, số lượng
thừa số \(m\) có thể tăng, nhưng số chữ số \(n\) của kết quả không đổi. Miễn
là dùng bộ đệm, số phép nhân vẫn xấp xỉ \(n(n-1)/2\). Vì vậy, riêng bước
phân tích thừa số nguyên tố không làm thuật toán chậm đi; điều này cũng phù
hợp với thực nghiệm.

Lợi ích của việc phân tích thừa số nguyên tố nằm ở chỗ sau đó xuất hiện rất
nhiều phép tính lũy thừa. Ta có thể tối ưu phần tính lũy thừa, và nhờ vậy
phương pháp phân tích thừa số nguyên tố mới phát huy hiệu quả.

Một quy trình tính \(N!\) theo hướng này là:
\begin{enumerate}
  \item Liệt kê các số nguyên tố \(p\le N\).
  \item Với mỗi \(p\), tính số mũ
  \[
    e_p=\sum_{i\ge 1}\left\lfloor\frac{N}{p^i}\right\rfloor.
  \]
  \item Tính các lũy thừa \(p^{e_p}\) bằng phương pháp lũy thừa nhị phân,
  dùng phép bình phương đã cải tiến.
  \item Điều chỉnh và ghép các lũy thừa khi có lợi, rồi nhân các khối kết
  quả bằng phép nhân độ chính xác cao có bộ đệm.
\end{enumerate}

\section{Lũy thừa nhị phân}

Lũy thừa nhị phân dùng các công thức
\[
  a^{2n+1}=a^{2n}\cdot a,\qquad
  a^{2n}=(a^n)^2,
\]
trong đó \(a\) là một số đơn chính xác.

Giả sử tích của một số \(n\) chữ số với một số \(m\) chữ số luôn có \(n+m\)
chữ số. Gọi \(F(n)\) là số phép nhân giữa chữ số cơ số cao cần dùng để tính
\(a^n\) bằng lũy thừa nhị phân. Khi đó
\[
  F(2n)=F(n)+n^2,\qquad F(1)=0.
\]
Nếu \(n=2^k\), ta có
\[
  F(n)=F(2^k)=\sum_{i=0}^{k-1}4^i
  =\frac{4^k-1}{3}
  =\frac{n^2-1}{3}.
\]

Trong khi đó, nhân liên tiếp cần \(n(n-1)/2\) phép nhân. Như vậy nhân liên
tiếp chỉ nhiều hơn lũy thừa nhị phân khoảng 50\%. Việc dùng lũy thừa nhị phân
không làm độ phức tạp giảm từ \(n^2\) xuống \(n\log n\), vì chi phí của các
phép bình phương số lớn vẫn chi phối.

\subsection{Cải tiến phép bình phương}

Ta khai thác đẳng thức
\[
  (a+b)^2=a^2+2ab+b^2.
\]
Chẳng hạn,
\[
  12345^2=123^2\cdot 10000+45^2+2\cdot 123\cdot 45\cdot 100.
\]
Ý tưởng là tách một số \(n\) chữ số thành một phần \(t\) chữ số và một phần
\(n-t\) chữ số, rồi tính bình phương bằng công thức trên.

Gọi \(F(n)\) là số phép nhân cần dùng để bình phương một số \(n\) chữ số. Ta
có
\[
  F(n)=F(t)+F(n-t)+t(n-t),\qquad F(1)=1.
\]
Bằng quy nạp toán học, suy ra
\[
  F(n)=\frac{n(n+1)}{2}.
\]
Vì thế, bất kể chia ở vị trí nào, hiệu quả về số phép nhân là như nhau. Trong
cài đặt, tách thành một phần 1 chữ số và một phần \(n-1\) chữ số là thuận
tiện hơn.

Cũng có thể nhìn phép bình phương theo công thức
\[
  S^2=
  \left(\sum_{0\le i<n} a_i10^i\right)^2
  =
  \sum_{0\le i<n}a_i^2 10^{2i}
  +2\sum_{0\le i<j<n}a_i a_j 10^{i+j}.
\]
Tổng cộng có \(n+\binom{n}{2}=n(n+1)/2\) phép nhân. Thuật toán bình thường
cần \(n^2\) phép nhân, tức chậm hơn khoảng một lần so với phép bình phương
đã cải tiến.

Nếu trong lũy thừa nhị phân ta dùng phép bình phương cải tiến, số phép nhân
cần dùng trở thành
\[
  \frac{(n+4)(n-1)}{6},
\]
xấp xỉ một phần ba so với thuật toán nhân liên tiếp với \(n(n-1)/2\) phép
nhân.

\section{Điều chỉnh sau khi phân tích thừa số nguyên tố}

\subsection{Vì sao cần điều chỉnh}

Khi tính
\[
  S=2^{11}3^{10},
\]
ta có thể tính riêng \(2^{11}\), tính riêng \(3^{10}\), rồi nhân hai kết quả.
Nhưng cũng có thể dùng
\[
  S=2^{11}3^{10}=6^{10}\cdot 2,
\]
tức tính \(6^{10}\) rồi nhân thêm \(2\). Hai cách này có hiệu quả khác nhau.

\subsection{Khi nào nên điều chỉnh}

Xét
\[
  S=a^{x+k}b^x=(ab)^x a^k.
\]
Khi
\[
  k < x\log_a b,
\]
dùng dạng \((ab)^x a^k\) tốt hơn; ngược lại, dùng dạng \(a^{x+k}b^x\) nhanh
hơn.

Có hai cách hiểu kết luận này. Cách trực tiếp là tính số phép nhân của hai
phương án rồi giải bất đẳng thức. Một cách nhìn khác là: khác biệt chính của
hai phương án nằm ở phép nhân cuối cùng. Vì tích cuối cùng của hai cách có
cùng số chữ số, hai thừa số càng chênh lệch thì số phép nhân cuối cùng càng
ít.

Cụ thể, ta chọn \((ab)^x a^k\) khi
\[
  a^x b^x-a^k > a^{x+k}-b^x.
\]
Bất đẳng thức này tương đương với
\[
  (b^x-a^k)(a^x+1)>0,
\]
nên tương đương với
\[
  b^x>a^k.
\]
Do đó điều kiện là
\[
  k < x\log_a b.
\]

\section{Kết luận}

Những điểm chính của bài viết là:
\begin{itemize}
  \item dùng \texttt{Comp} làm kiểu cơ sở cho mỗi đoạn nhỏ;
  \item cân nhắc biểu diễn nhị phân, dù trong thi đấu thường phải trả giá ở
  bước đổi sang hệ thập phân;
  \item dùng bộ đệm khi nhân liên tiếp độ chính xác cao, và chọn kích thước
  bộ đệm bằng kích thước mỗi đoạn nhỏ;
  \item dùng phép bình phương cải tiến;
  \item dùng lũy thừa nhị phân;
  \item tính giai thừa bằng phân tích thừa số nguyên tố, kèm điều chỉnh sau
  khi phân tích.
\end{itemize}

Các kỹ thuật này có thể dùng cho tính lũy thừa độ chính xác cao, tính tích
liên tiếp như giai thừa, và tính số chỉnh hợp hoặc tổ hợp độ chính xác cao.

Tự thân bài toán tính \(N!\) bằng độ chính xác cao không khó, nhưng vẫn có
thể được tối ưu từ nhiều góc độ. Thật ra, nhiều thuật toán kinh điển cũng còn
không gian cải thiện, và những chương trình do ta tự viết cũng vậy. Chỉ cần
chịu khó phân tích, rất có thể ta sẽ tìm ra một thuật toán tốt hơn. Trong thi
đấu, yêu cầu về độ chính xác cao thường không quá nặng, nên các tối ưu này có
thể không bắt buộc; mục tiêu chính của ví dụ này là nhấn mạnh một ý: thuật
toán luôn có thể được tối ưu.

\end{document}
